% from ega4-18-fr.tex
\oldpage[IV-4]{156}

l'ensemble sous-jacent à $f^{-1}(\eta)$, et la fibre $f^{-1}(\eta)$ est somme des $\Spec(k(\xi_i))$. L'hypothèse sur les $X_y$ et le fait que ces fibres sont de dimension 0, montrent qu'elles sont \emph{géométriquement régulières} (6.7.6), et \emph{a fortiori} géométriquement normales, autrement dit $f$ est un \emph{morphisme normal}. Mais puisque $A$ est intègre et géométriquement unibranche, il résulte alors de (6.15.10) que $B$ est intègre et géométriquement unibranche, donc $f^{-1}(\eta)=X-\{b\}$ est réduit à un seul point, et \emph{a fortiori} connexe; ceci termine donc la démonstration de (18.9.7.3) et celle de (18.9.7).

\begin{remark}[18.9.7.8]
\label{IV.18.9.7.8}
Dans le cas C) de la démonstration de (18.9.7.3), on peut éviter de faire appel au délicat résultat (6.15.10) en raisonnant de la façon suivante: puisque $A$ est intègre et de dimension 1, il résulte du théorème de Krull-Akizuki (Bourbaki, \emph{Alg. comm.}, chap.~VII, \textsection~2, n\textsuperscript{o}~5, prop.~5) que sa clôture intégrale $A'$ est un anneau \emph{noethérien}. Le même raisonnement qu'au début de la démonstration de (18.9.7) montre alors que $B'=B\otimes_A A'$ est un anneau local; en outre, par platitude, $B'$ est contenu dans l'anneau total des fractions $R$ de l'anneau réduit $B$ (3.3.5). Or, le raisonnement qui prouve le théorème de Krull-Akizuki (Bourbaki, \emph{loc. cit.}) s'applique également à un anneau local noethérien \emph{réduit} de dimension 1, et montre que pour un tel anneau $B$, tout anneau compris entre $B$ et son anneau total des fractions est \emph{noethérien}. L'anneau $B'$ étant un anneau local noethérien et le morphisme $\Spec(B')\to\Spec(A')$ étant normal, $B'$ est un anneau intègre et intégralement clos (6.5.4), et on conclut comme au début de la preuve de (18.9.7).
\end{remark}

\begin{corollary}[18.9.8]
\label{IV.18.9.8}
Avec les notations de (18.9.7), supposons vérifiée la condition (i) de (18.9.7) et l'une des deux conditions suivantes:
\begin{enumerate}
  \item[(ii$'$)] $A$ est un anneau strictement local (18.8.2).
  \item[(ii$''$)] $A$ est un anneau local hensélien et $k(B)$ est une extension primaire de $k(A)$.
\end{enumerate}
Alors toutes les fibres $X_y$ ($y\in Y$) sont géométriquement connexes.
\end{corollary}

Soit $\mathfrak p$ l'idéal premier de $A$ correspondant au point $y$, et posons $A'=A/\mathfrak p$ et $B'=B\otimes_A A'=B/\mathfrak pB$, qui sont évidemment des anneaux locaux noethériens; il est clair que le morphisme $f':\Spec(B')\to\Spec(A')$ vérifie la condition (i) de (18.9.7); comme $y$ est le point générique de $\Spec(A')$ et que $k(A')=k(A)$, $k(B')=k(B)$, on voit qu'il suffit de vérifier que, dans le cas (ii$'$) (resp. (ii$''$)), $A'$ est géométriquement unibranche (resp. unibranche). Mais $A'$ est intègre, et dans le cas (ii$'$) (resp. (ii$''$)) il est strictement local (resp. hensélien) en vertu de (18.5.10). On en conclut que $A'$ est géométriquement unibranche (resp. unibranche) en vertu de (18.8.15) (resp. (18.6.12)).

\begin{corollary}[18.9.9]
\label{IV.18.9.9}
Soit $A$ un anneau local noethérien hensélien, universellement japonais (\emph{i.e.} ((7.6.4) et (7.7.2)), dont les fibres formelles (7.3.13) sont géométriquement réduites); alors les fibres formelles de $A$ sont géométriquement connexes.
\end{corollary}

Il suffit d'appliquer (18.9.8) au cas où $B=\widehat A$, puisque $B$ est un $A$-module plat et que $k(B)=k(A)$.

\begin{remark}[18.9.10]
\label{IV.18.9.10}
Les démonstrations de (18.9.4) et (18.9.5) n'utilisent que le fait que les fibres formelles de $A$ sont géométriquement connexes. Les conclusions

